% Generated aggregate fragment; not yet included in the manuscript.
\paragraph{Fixed-input event-structure diagnostic.}
A retrospective NYC audit fixes four core and twelve buffer rows,
exact public times and query values, and compares ordered events,
exactly-two-row events, and clique events at $C\in\{2,3,4\}$.
Only positive buffer counts feasible in all three families are used
for outcome comparisons, with the same count on each side.
The audit certifies 36 endpoint pairs;
0 of 24 fixed-count comparisons change an endpoint.
\begin{table}[t]
\centering
\begin{tabular}{c|rrr}
$C$ & Ordered & Two-row & Clique \\
\hline
2 & 4 & 4 & 4 \\
3 & 8 & 4 & 8 \\
4 & 12 & 4 & 12 \\
\end{tabular}
\caption{Maximum selected-buffer counts on one fixed NYC small audit input. These model-specific maxima are not fixed-count outcome comparisons.}
\label{tab:nyc-structure-diagnostic}
\end{table}
All sixteen intervals share 14 seconds of common overlap.
Thus every admissible ordered event is already a clique; this
instance does not test the value of connected non-clique events.
The result establishes neither a distinct public-data outcome
advantage for ordered structure nor a general equivalence of the models.
Buffer selection requires complete public miles and trip-time values.
No true-event recovery, true-world coverage or population claim follows.
